% 双阈值比较的教学序列：第10帧硬切，第22至31帧叠化；所有差值均为构造数值。宽约100mm。
\begin{tikzpicture}[font=\figurefont\small]
  \begin{axis}[
    width=100mm,height=48mm,
    xmin=0,xmax=41,ymin=0,ymax=1.1,
    xlabel={帧序号$t$},ylabel={差值},
    xtick={1,10,22,31,40},ytick={0.08,0.5,1},
    yticklabels={$T_{\mathrm l}$,$T_{\mathrm h}$,1},
    legend cell align=left,
    legend style={font=\figurefont\footnotesize,at={(0.5,-0.36)},anchor=north,
      legend columns=2,draw=none}]
    \addplot[ybar,area legend,bar width=1.2mm,draw=FigureBlue,fill=FigureBlueFill] coordinates {
      (1,0.04)(2,0.05)(3,0.03)(4,0.06)(5,0.04)(6,0.05)(7,0.04)(8,0.03)(9,0.05)
      (10,1.0)
      (11,0.04)(12,0.05)(13,0.06)(14,0.04)(15,0.03)(16,0.05)(17,0.04)(18,0.05)
      (19,0.04)(20,0.06)(21,0.05)
      (22,0.11)(23,0.12)(24,0.10)(25,0.12)(26,0.11)(27,0.12)(28,0.10)(29,0.11)
      (30,0.12)(31,0.11)
      (32,0.05)(33,0.04)(34,0.05)(35,0.03)(36,0.04)(37,0.05)(38,0.04)(39,0.06)
      (40,0.04)};
    \addlegendentry{相邻帧差$d(t)$}
    \addplot[FigureRed,line width=1.1pt,mark=*,mark size=1.2pt] coordinates {
      (22,0.11)(23,0.22)(24,0.32)(25,0.42)(26,0.52)(27,0.61)(28,0.70)(29,0.78)
      (30,0.86)(31,0.93)};
    \addlegendentry{与起点第21帧的差}
    \addplot[FigureMuted,dashed,line width=.6pt,domain=0:41] {0.5};
    \addplot[FigureMuted,dashed,line width=.6pt,domain=0:41] {0.08};
  \end{axis}
\end{tikzpicture}
